\section{Power Amplifiers Output Stages / Classes}

When consulting amplifier datasheets, one of the first pieces of information given is the \textbf{amplifier class}: Class A, B, AB, etc.  
\begin{figure}[H]
    \centering
    \includegraphics[width=0.7\linewidth]{Signal_Chain.png}
    \caption{Signal chain.}
\end{figure} 
\[
\text{MIC} \;\rightarrow\; \text{PRE} \;\rightarrow\; \text{FILTER} \;\rightarrow\; \text{EQ} \;\rightarrow\; \text{POST-EQ PROCESSING} \;\rightarrow\; \text{POWER AMP} \;\rightarrow\; \text{LOUDSPEAKER}
\]
\textbf{Note:}  
The power amplifier is usually configured as a \textbf{voltage buffer} and is mainly responsible for \textbf{current amplification}.

\subsection{Class A Speaker}

\begin{figure}[H]
    \centering
    % Left side: explanation + equations
    \begin{minipage}[c]{0.65 \linewidth}
    This amplifier class is based on an \textbf{emitter-follower architecture}.  
    \begin{itemize}
        \item $R_L$ is the load impedance (the loudspeaker).
        \item The capacitor $C$ is used to decouple the DC component, so that the loudspeaker only receives the AC part of the signal (average level $\approx 0$).
    \end{itemize}
    \medskip
    If the capacitor $C$ is large enough, the output node satisfies:
    \[
    v_{out} \approx v_{E} \quad \text{(AC behavior)}
    \]
    \end{minipage}\hfill
    \begin{minipage}[c]{0.25\linewidth}
        \centering
        \includegraphics[width=\linewidth]{Class_A_Stage.png}
    \caption{Circuit of the class A stage.}
    \end{minipage}\hfill
\end{figure}

Indeed, we have: 
\[
v_{out} = v_E \frac{R_L}{R_L+1/(sC)}
\quad \Rightarrow \quad
\frac{v_{out}}{v_E} = \frac{sCR_L}{1+sCR_L}  \ \underset{C\rightarrow+\infty}{\rightarrow} \  1
\]
Moreover, we know that:
\[
v_{out} = R_L \cdot  i_L \approx R_L \cdot i_C, \quad \text{assuming } i_Q\approx 0
\]
Since $ i_C = g_m v_{BE} $, we obtain 
\begin{align*}
    & v_{out} = R_L \cdot g_m \cdot v_{BE} = R_L \cdot g_m \cdot \left( v_B - v_E\right) 
\quad \Rightarrow \quad 
v_{out} = R_L \cdot g_m \cdot \left( v_{in} - v_{out}\right) \quad \text{since } v_E \approx v_{out}  \\
   \Rightarrow \quad & v_{out} \left( 1+ g_mR_L\right) = R_L \cdot g_m \cdot v_{in} 
   \quad \Rightarrow \quad
   \frac{v_{out}}{v_{in}} = \frac{g_m \cdot R_L}{1+g_m \cdot R_L} = A_V
\end{align*}
The \textbf{voltage transfer} is:
\[
\frac{v_{out}}{v_{in}} = \frac{g_m \cdot R_L}{1 + g_m \cdot R_L} \approx 1
\]
since we assume $g_m \cdot R_L \gg 1$.  
\textit{(This is similar to the principle of the tube screamer circuit.)}

The \textbf{Current Gain} is: 
\[
i_L = i_E = (1 + \beta) i_B
\]
since the quiescent current $I_Q$ is disconnected in this configuration.  

Thus:
\[
\frac{i_E}{i_B} = \beta + 1 = A_I \gg 1
\]
In practice, $\beta + 1$ typically ranges between $10$ and $100$.  
This shows that a Class A stage provides \textbf{large current gain} but only \textbf{unity voltage gain}.

\paragraph{Computing the Required Voltage Supply:}

Let us consider a power amplifier designed to deliver the following load power:
\[
P_L = 100 \, W \quad \text{with} \quad R_L = 8 \, \Omega
\]
The load power is related to the RMS voltage:
    \[
    P_L = \frac{v_{RMS,L}^2}{R_L}
    \quad \Rightarrow \quad v_{RMS,L} = \sqrt{P_L \cdot R_L}
    \]
    Substituting values:
    \[
    v_{RMS,L} = \sqrt{100 \cdot 8} = 28.3 \, V_{RMS}
    \]
\begin{figure}[H]
    \centering
    % Left side: explanation + equations
    \begin{minipage}[c]{0.65 \linewidth}
    
    The corresponding peak voltage is:
    \[
    \hat{V}_L = \sqrt{2} \cdot v_{RMS,L} = 40 \, V
    \]
    Therefore, the amplifier must provide at least $40 \, V$ peak at the output.  
    As a supply condition, the total supply voltage must satisfy:
    \[
    V_{CC} \geq 2 \cdot \hat{V}_L = 80 \, V
    \]
    \end{minipage}\hfill
    \begin{minipage}[c]{0.3\linewidth}
        \centering
        \includegraphics[width=\linewidth]{V_CC_Class_A.png}
    \caption{$V_{CC}$ representation.}
    \end{minipage}\hfill
\end{figure}

\paragraph{Computing the Required Current Supply:}

Using the same configuration, the load current is:
\[
\hat{i}_L = \frac{\hat{V}_L}{R_L} = \frac{40}{8} = 5 \, A
\]
Thus, the power supply must be capable of delivering a \textbf{peak current of 5 A}.

Finally: 
\[
\hat{i}_{IN} = \frac{\hat{i}_L}{\beta + 1} \approx 50 \ \text{mA assuming that} \ \beta + 1 =100
\]

\subsection{Efficiency of an Output Stage}

The \textbf{efficiency} of a power amplifier is defined as:
\[
\eta \in [0,1], \qquad \eta = \frac{P_L}{P_S}
\]
where:
\begin{itemize}
    \item $P_L$ = average power delivered to the load (output power)
    \item $P_S$ = average power drawn from the supply (input power)
\end{itemize}

For a sinusoidal signal, the \textbf{load power} is:
\[
P_L = \frac{v_{RMS,L}^2}{R_L} = \frac{\hat{V}_L^2}{2R_L}
\]
At maximum efficiency, the peak load voltage satisfies:
\[
\hat{V}_L = \frac{V_{CC}}{2}
\quad \Rightarrow \quad
P_L = \frac{V_{CC}^2}{8R_L}
\]
The \textbf{average power delivered by the supply} is:
\[
P_S = \frac{1}{T} \int_0^T V_{CC} \cdot i_S \, dt 
= V_{CC} \cdot \frac{1}{T} \int_0^T i_S \, dt
= V_{CC} \cdot \overline{I_S}
\]

where $\overline{I_S}$ is the average current drawn from the supply.

We can think of $i_S$ as the superposition of:
\[
i_S = i_L + i_Q
\]
where: 
\begin{itemize}
    \item $i_L$: the load current (sinusoidal)  
    \item $i_Q$: the quiescent current of the transistor \textbf{assumed to be constant}
\end{itemize}

\begin{figure}[H]
    \centering
    % Left side: explanation + equations
    \begin{minipage}[c]{0.65 \linewidth}
    As $i_L$ is sinusoidal, its average is zero:
    \[
    \overline{I_L} = 0
    \quad \Rightarrow \quad 
    \overline{I_S} = I_Q
    \]
    Since t\textbf{he supply current and the transistor current cannot be negative}, the load must provide current to $i_Q$ when the transistor is off. I.e.
    \[
    \check{I}_S \geq 0 
    \quad \Rightarrow \quad
    \check{I}_L + \check{I}_Q = -\frac{V_{CC}}{2R_L}+\check{I}_Q \geq 0 
    \quad \Rightarrow \quad  
    \check{I}_Q \geq \frac{V_{CC}}{2 R_L}
    \]
    Thus minimum quiescent current is:
    \[
    I_Q^{min} = \frac{V_{CC}}{2R_L}
    \]
    For the peak value, we have:
    \[
    \hat{I}_S = I_Q + \hat{i}_L = I_Q + \frac{\hat{v}_L}{R_L} = I_Q + \frac{{V}_{CC}}{2 R_L}
    \]
    \end{minipage}\hfill
    \begin{minipage}[c]{0.3\linewidth}
        \centering
        \includegraphics[width=\linewidth]{I_s(t).png}
    \caption{Graph of $I_s(t)$.}
    \end{minipage}\hfill
\end{figure}
Note: $\check{I}$ denotes the minimum value of $I$.

Thus: 
    \[{{P}_{s}}={{V}_{CC}}\cdot {{\bar{I}}_{s}}={{V}_{CC}}\frac{{{V}_{CC}}}{2{{R}_{L}}}=\frac{V_{CC}^{2}}{2{{R}_{L}}}
    \quad \Rightarrow \quad {\eta }_{max}=\frac{V_{CC}^{2}}{8{{R}_{L}}}\frac{2{{R}_{L}}}{V_{CC}^{2}}=\frac{1}{4}
    \]
The power is wasted since the load sometimes needs a negative current while we keep feeding a positive current (\textbf{TO
BE RECHECKED}). To solve this matter, we introduced class B stages.

\subsection{Class B Stages}

From the previous analysis, the efficiency of a Class A stage is limited to:
\[
\eta_{max} = \frac{1}{4} = 25\%
\]
This low efficiency arises because the load sometimes requires negative current while the supply continues delivering positive current, leading to wasted power.  
To overcome this issue, \textbf{Class B stages} were introduced.

\begin{figure}[H]
    \centering
    % Left side: explanation + equations
    \begin{minipage}[c]{0.65 \linewidth}
    \textbf{Principle of Operation:}
Class B stages are designed to provide much higher efficiency compared to Class A.  
\begin{itemize}
    \item Instead of a single amplifier, Class B uses a \textbf{push-pull configuration}.
    \item The circuit employs two complementary transistors:
    \begin{itemize}
    \item An NPN transistor (top) that conducts during the positive half-cycle.  
    \item A PNP transistor (bottom) that conducts during the negative half-cycle.  
  \end{itemize}
\end{itemize}
This way, each transistor only amplifies one half of the input waveform.
    \end{minipage}\hfill
    \begin{minipage}[c]{0.3\linewidth}
        \centering
        \includegraphics[width=\linewidth]{Class_B.png}
    \caption{Circuit of the B stage.}
    \end{minipage}\hfill
\end{figure}


\begin{figure}[H]
    \centering
    % Left side: explanation + equations
    \begin{minipage}[c]{0.3 \linewidth}
    \centering
        \includegraphics[width=\linewidth]{class_B_graph.png}
    \caption{Graph of $v_{in},v_{out}$ for the class B stage.}
    \end{minipage}\hfill
    \begin{minipage}[c]{0.65\linewidth}
        \textbf{Conduction Conditions:}
    \[
    \begin{cases}
    Q_1 \ \text{(NPN) is ON if } V_{BE} \geq 0.7 \, V \\
    Q_2 \ \text{(PNP) is ON if } V_{EB} \geq 0.7 \, V \ \Leftrightarrow \ V_{BE} \leq -0.7 \, V
    \end{cases}
    \]
    \end{minipage}\hfill
\end{figure}

\textbf{Crossover Distortion:}
At low input voltages (close to $0$), neither transistor is conducting.  
This creates a \textbf{dead zone} around the origin, which causes \textbf{crossover distortion}.  

Crossover distortion is a common drawback of pure Class B stages, leading to signal non-linearity when the input amplitude is small. 

To correct it, we introduced AB stages.

\subsection{Class AB speaker}

\begin{figure}[H]
    \centering
    % Left side: Text
    \begin{minipage}[c]{0.44\textwidth}
    \textbf{Class AB Solution:}
To avoid crossover distortion, we can apply a small \textbf{DC bias shift} to the transistor bases (about 0.7 V).  
        This ensures that both transistors conduct slightly even when the input is near zero, thus eliminating the dead zone.  

        The resulting stage is called a \textbf{Class AB amplifier}, which combines:
        \begin{itemize}
            \item The higher efficiency of Class B,  
            \item With the low distortion of Class A.
        \end{itemize}
    \end{minipage}
    \hfill
    % Right side: Image
    \begin{minipage}[c]{0.55\textwidth}
        \centering
        \includegraphics[width=\linewidth]{class AB power amplifier.png}
        \caption{Scheme of the AB class (left) and graph of $v_{in}, v_{out}$ (right)}
        \label{fig:placeholder}
    \end{minipage}
\end{figure}
The \textbf{Class AB configuration} is essentially a refinement of Class B.  
It removes most of the crossover distortion by biasing both transistors slightly on, although small artifacts may still remain.  
This stage uses a \textbf{push-pull configuration}: one transistor sources current (push), while the other sinks current (pull).

The efficiency is:
\[
\eta = \frac{P_L}{P_S} 
= \frac{V_{CC}^2 / (8 R_L)}{\frac{1}{T} \int_0^T V_{s} \cdot I_S \, dt}
= \frac{{V}_{CC}^2 / (8R_L)}{V_{CC} \cdot \overline{I_S}}
\quad \Rightarrow \quad 
\eta = \frac{{V}_{CC}}{8 R_L \overline{I_S}}
\]
For a sinusoidal waveform (positive half-wave only), the average supply current is:
\[
\overline{I_S} = \frac{1}{T} \int_0^{T/2} I_S \, dt
=
\frac{1}{T} \int_0^{T/2} \frac{V_{CC}}{2R_L}\sin(\omega t) \, dt
\quad \Rightarrow \quad
\overline{I_S}= \frac{V_{CC}}{2R_L} \cdot \frac{1}{T} \int_0^{T/2} \sin(\omega t) \, dt, \qquad \omega = \frac{2 \pi}{T}
\]
Evaluating the integral:
\[
\overline{I_S} = \frac{V_{CC}}{2 \cdot R_L \cdot T} \left[-\frac{\cos(\omega t)}{\omega} \right]_0^{T/2}
= \frac{V_{CC}}{2R_L} \frac{1}{T \tfrac{2 \pi}{T}}\cdot [1 + 1]
= \frac{V_{CC}}{2 R_L \pi}
\]
Therefore, the power is : 
\[
P_S = V_{CC} \cdot \overline{I_S} = \frac{V_{CC}^2}{2 R_L \pi}
\]
The efficiency is obtained through: 
\[
\eta = \frac{P_L}{P_S}
= \frac{\frac{V_{CC}^2}{8R_L}}{\frac{V_{CC}^2}{2 \pi R_L }}
= \frac{2 \pi R_L}{8 R_L} = \frac{\pi}{4} \approx 79 \%
\]
Note that for the given conditions, we have:
\[
P_{supply}|_{P_L=100W} \approx 127 \, W
\]
This shows that Class AB is much more efficient than Class A (25\%), while keeping distortion at a low and acceptable level.

\textbf{Remark:}
Class AB stages may still be subject to \textbf{thermal runaway}, requiring bias stabilization mechanisms.


\subsection{Class D Stages}

Class D amplifiers use a more complex architecture, but they rely on two fundamental principles:
\begin{itemize}
    \item Very efficient ``digital'' voltage transfer ($V_{CC}$ or $0V$).  
    \item Very efficient passive low-pass filters (LC filters rather than RC, since we want to avoid power dissipation).  
\end{itemize}

\textbf{Example:}
Consider the simple circuit below, where we want:
$v_{out} = V_{CC}$. By introducing a switch between the supply and the output, we can directly connect or disconnect the output from the supply.  
For high efficiency, the switch must have an extremely small resistance when closed.  

In the ideal case:
\[
\eta = \frac{P_L}{P_S} \approx 1 \quad \Rightarrow \quad \eta = 100\%
\]

\begin{figure}[H]
    \centering
    \includegraphics[width=0.35\linewidth]{Class_D.png}
    \caption{Circuit of Class D stage.}
    \label{fig:placeholder}
\end{figure}

\clearpage

\subsubsection{Pulse Width Modulation (PWM)}
To achieve a ``digital'' representation of $v_{out}$, Class D amplifiers use \textbf{Pulse Width Modulation (PWM)}.
\begin{figure}[H]
    \centering
    % Left side: explanation + equations
    \begin{minipage}[c]{0.55 \linewidth}
The duty cycle $\alpha$ is defined as:
\[
\alpha = \frac{T_H}{T_H + T_L} = \frac{T_H}{T}
\]
where $T_H$ is the time the switch is ON, and $T_L$ the time it is OFF.  

Efficiency is maximum when the PWM is $0$ (output OFF), and minimum when it is fully ON.  
    \end{minipage}\hfill
    \begin{minipage}[c]{0.4\linewidth}
        \centering
        \includegraphics[width=\linewidth]{PWM_graph.png}
    \caption{Graph of a PWM.}
    \end{minipage}\hfill
\end{figure}
The average value of the PWM output is:
\[
\overline{v_{out}} = V_{CC} \cdot \alpha
\]

\textbf{Principle of Modulation:}
The idea is to vary the duty cycle in order to encode a time-varying signal.  
\begin{itemize}
    \item For instance, a sinusoidal input can be reconstructed from the average value of a PWM signal whose duty cycle follows the sinusoid. 
    \item The PWM is typically generated by comparing the input signal with a \textbf{carrier wave} (usually triangular). 
    \item The carrier signal $v_{carrier}$ has a very high frequency (in the MHz range), much higher than the input audio signal.
\end{itemize}

\begin{figure}[H]
    \centering
    % Left side: explanation + equations
    \begin{minipage}[c]{0.45 \linewidth}
        \centering        \includegraphics[width=\linewidth]{PWM_1.png}
    \caption{Comparison with rectangular carrier waves.}  
    \end{minipage}\hfill
    \begin{minipage}[c]{0.48\linewidth}
        \centering
        \includegraphics[width=\linewidth]{PWM_2.png}
    \caption{Comparison with triangular carrier waves.}
    \end{minipage}\hfill
\end{figure}



\textbf{Filtering and Output:}
The PWM signal is then passed through a \textbf{low-pass filter} (LC filter) that extracts the average value.  
The result is an amplified version of the input signal, with very high efficiency and low power dissipation in the amplifier itself.

\begin{figure}[H]
    \centering
    \includegraphics[width=0.8\linewidth]{PWM_Pipeline.png}
    \caption{Overall PWM pipeline.}
    \label{fig:placeholder}
\end{figure}

\subsection*{From the Power Amplifiers Slides}

The load voltage can be expressed as:
\[
v_L = \frac{R_L}{R_0 + R_L} \cdot v_0
\]
where:
\begin{itemize}
    \item $R_0$ is the internal (output) resistance of the amplifier,  
    \item $R_L$ is the load resistance (the loudspeaker),  
    \item $v_0$ is the internal voltage source of the amplifier.
\end{itemize}

\begin{figure}[H]
    \centering
    % Left side: explanation + equations
    \begin{minipage}[c]{0.45 \linewidth}
        
        \centering
        \includegraphics[width=\linewidth]{P(R_L).png}
    \caption{Graph of $P(R_L)$.}
    \end{minipage}\hfill
    \begin{minipage}[c]{0.47\linewidth}
    The load power is given by:
        \[
        P_L = v_L \cdot i_L
        \]
        Since:
        \[
        i_L = \frac{v_L}{R_L} = \frac{v_0}{R_0 + R_L}
        \]
        we can write:
        \[
        P_L = \frac{v_0^2 R_L}{(R_0 + R_L)^2}
        \]
        From this expression, the maximum power delivered to the load (\textbf{maximum power transfer}) occurs when:
        \[
        R_L = R_0
        \]
        At this point, the power transfer from the amplifier to the load is optimal.
    \end{minipage}\hfill
\end{figure}


\textbf{Practical Limitation:}
In practice, achieving $R_L = R_0$ across the entire audio bandwidth is \textbf{impossible}, because:
\begin{itemize}
    \item The load resistance $R_L$ (i.e., the loudspeaker impedance) is not constant but varies strongly with frequency.  
    \item This mismatch means that the condition for maximum power transfer can only be satisfied at a specific frequency, not across the whole spectrum.
\end{itemize}


\textbf{Key Observation:}
\begin{itemize}
    \item If $R_L \gg R_0$: most voltage appears across the load, but the power transfer is low.  
    \item If $R_L \ll R_0$: most power is dissipated inside the amplifier, reducing efficiency.
    \item Only at $R_L = R_0$ do we reach maximum transfer, which is why amplifier–loudspeaker matching is a key design consideration in audio systems.
\end{itemize} 

\subsection{Notes from the slides}
Power amplifiers are the final stage in an audio signal chain.  
Their purpose is to deliver sufficient voltage and current to drive loudspeakers, ensuring accurate reproduction of the input signal at the desired sound pressure level.  
While small-signal amplifiers focus on linearity and low noise, power stages must also handle high power transfer, thermal dissipation, and protection against overloads.


\paragraph{Example: Luxman 900u Integrated Amplifier}
The \textbf{Luxman 900u} is a high-end reference amplifier often cited in audio engineering for its design precision and sonic transparency.  
It exemplifies the principles of Class AB topology optimized for low distortion and wide bandwidth.

\begin{itemize}
  \item \textbf{Topology:} Dual-mono symmetrical configuration to minimize crosstalk and maintain channel isolation.
  \item \textbf{Amplification stage:} Parallel complementary push-pull output pairs using high-speed bipolar transistors.
  \item \textbf{Biasing:} Precisely controlled quiescent current ensures Class A operation at low levels, transitioning to Class AB at higher output.
  \item \textbf{Feedback:} Local and global feedback networks maintain low total harmonic distortion (THD < 0.005\%) across the full audio band.
  \item \textbf{Power supply:} Oversized toroidal transformers with independent filtering for each channel reduce hum and noise coupling.
\end{itemize}


\paragraph{Key Performance Characteristics.}
\begin{itemize}
  \item Power output: typically $120\,\text{W}/\text{channel}$ at $8\,\Omega$.
  \item Frequency response: $10\,\text{Hz}$ – $100\,\text{kHz}$ (–1 dB).
  \item Damping factor: above 150, ensuring tight control of loudspeaker motion.
  \item Slew rate: high (tens of V/µs), allowing fast transient response.
\end{itemize}

\paragraph{Output Protection Mechanisms: }
High-power amplifiers must protect both the output devices and the load (speakers) from electrical faults or excessive currents.  
Typical \textbf{short-circuit protection} strategies include:
\begin{itemize}
  \item \textbf{Current sensing resistors} monitoring the emitter current in each output transistor.
  \item \textbf{Comparator and limiter circuits} that reduce drive or cut the signal when current exceeds a safe threshold.
  \item \textbf{Output relays} that disconnect the load under DC offset or startup transients.
\end{itemize}

\paragraph{Power Transistors and Output Devices}
Modern power amplifiers use complementary NPN–PNP (or N-channel–P-channel MOSFET) pairs for the final stage.  
Desirable properties of power devices include:
\begin{itemize}
  \item High current handling and safe operating area (SOA),
  \item High transition frequency ($f_T$) for wide bandwidth,
  \item Low output capacitance and minimal storage time (for low distortion),
  \item Thermal stability via bias-compensating networks.
\end{itemize}

\paragraph{BJT vs. MOSFET Output Stages.}
\begin{itemize}
  \item \textbf{BJTs:} offer lower on-resistance and better linearity; require bias stabilization to avoid thermal runaway.
  \item \textbf{MOSFETs:} have voltage-controlled gates and inherent thermal stability, making them suitable for high-speed or Class D designs.
\end{itemize}

\paragraph{Power Transfer and Load Matching}
The load impedance strongly influences power delivery.  
For an amplifier with open-circuit voltage $V_o$ and internal resistance $R_o$, the power transferred to a load $R_L$ is:
\[
P_{L} = \frac{V_o^2 R_L}{(R_L + R_o)^2}.
\]

\begin{figure}[H]
\centering
\includegraphics[width=0.4\textwidth]{Power_Load_Load_Resistance.png}
\caption{Example of power delivered to the load as a function of load resistance for $V_o = 20\,\text{V}$ and $R_o = 0.5\,\Omega$.  
Maximum power occurs when $R_L = R_o$, but in audio amplifiers $R_o \ll R_L$ to minimize losses and distortion.}
\end{figure}

In practice, power amplifiers are designed with a very low output impedance (high damping factor) to maximize voltage transfer rather than power matching, ensuring accurate loudspeaker control.

\paragraph{Class D Amplifiers and Applications:}
\textbf{Class D amplifiers} achieve high efficiency by using pulse-width modulation (PWM) and switching output devices rather than linear conduction.  
They are widely used in modern sound reinforcement systems and active loudspeakers.

\textbf{Key Characteristics.}
\begin{itemize}
  \item Efficiency above 85–90\%, drastically reducing heat generation.
  \item Compact size and low weight, ideal for portable and embedded systems.
  \item Require output low-pass filters to remove high-frequency switching components.
\end{itemize}

\begin{figure}[H]
\centering
\includegraphics[width=0.4\textwidth]{Class_D_System.png}
\caption{Simplified diagram of a Class D amplifier system driving a loudspeaker.  
The PWM signal is demodulated by the output LC filter, reconstructing the amplified audio waveform.}
\end{figure}


\paragraph{Applications in Professional Audio: }
Manufacturers such as \textbf{K-Array} and other live-sound brands implement Class D amplifier modules in powered speaker arrays and subwoofer drivers, achieving high output power with minimal weight and heat dissipation.

\begin{itemize}
  \item \textbf{Loudspeaker driver amplifiers:} each driver (woofer, mid, tweeter) is powered by its dedicated amplifier channel, allowing precise crossover and dynamic control.
  \item \textbf{Subwoofer amplifiers:} designed for very low frequencies, typically using high-current Class D stages capable of several hundred watts.
  \item \textbf{Integrated systems:} include DSP crossovers, limiters, and monitoring networks.
\end{itemize}


\textbf{Summary and Comparative Overview}
\begin{itemize}
  \item The Luxman 900u illustrates state-of-the-art linear amplification with extremely low THD and wide bandwidth.  
  \item Protection circuits (short-circuit, DC offset, over-temperature) are essential for reliability in high-power designs.
  \item The choice of output devices (BJT or MOSFET) affects both sound character and thermal behavior.
  \item Class D amplifiers represent the modern trend in professional audio for efficiency and integration, particularly in powered speaker systems.
  \item Proper load management and high damping factors ensure faithful control of loudspeaker motion and accurate reproduction of transient dynamics.
\end{itemize}
